\documentclass[11pt]{article}

%------------------------------------------------
% Pacotes
%------------------------------------------------
\usepackage[utf8]{inputenc}
\usepackage[T1]{fontenc}
\usepackage{lmodern}
\usepackage[brazilian]{babel}
\usepackage{amsmath,amssymb,amsthm,mathtools}
\usepackage{geometry}
\usepackage{microtype}
\usepackage{enumitem}
\usepackage{booktabs}
\usepackage{hyperref}

\geometry{margin=2.6cm}

%------------------------------------------------
% Ambientes
%------------------------------------------------
\theoremstyle{plain}
\newtheorem{theorem}{Teorema}[section]
\newtheorem{proposition}[theorem]{Proposição}
\newtheorem{lemma}[theorem]{Lema}
\newtheorem{corollary}[theorem]{Corolário}
\newtheorem{conjecture}[theorem]{Conjectura}
\newtheorem*{conjLY}{Conjectura de Lijun--Yun}
\theoremstyle{definition}
\newtheorem{definition}[theorem]{Definição}
\newtheorem{example}[theorem]{Exemplo}
\newtheorem{problem}[theorem]{Problema}
\theoremstyle{remark}
\newtheorem{remark}[theorem]{Observação}

%------------------------------------------------
% Comandos
%------------------------------------------------
\newcommand{\R}{\mathbb{R}}
\newcommand{\C}{\mathbb{C}}
\newcommand{\Q}{\mathbb{Q}}
\newcommand{\fe}{f_{\mathrm e}}
\newcommand{\fo}{f_{\mathrm o}}
\newcommand{\ip}[2]{\left\langle #1,#2\right\rangle}
\newcommand{\sgn}{\operatorname{sgn}}

\title{Generalizações de um critério de centro para equações de Abel:\\
momentos, ortogonalidade ponderada e condições de ordem superior\\[2mm]
\large Notas de trabalho a partir de Alves--Lemos--Araujo (2020)}
\author{Anderson L. A. de Araujo}
\date{Outubro de 2026}

\begin{document}

\maketitle

\begin{abstract}
Organizamos as possíveis generalizações do artigo de Alves, Lemos e Araujo
\cite{ALA20}, que dá uma resposta parcial positiva a uma conjectura de Lijun e
Yun \cite{LY01} sobre a equação de Abel
$x'=f(t)x^3+t^{n-1}x^2$, $t\in[-1,1]$. Partimos de duas extensões já
redigidas: (i) parte par de $f$ com grau $\le 4$ e parte ímpar arbitrária;
(ii) parte par com grau $\le 6$ sob uma condição de ortogonalidade ponderada.
Mostramos que ambas se encaixam num único esquema e o estendemos em quatro
direções:
(a) fórmulas fechadas para os coeficientes de centro de ordens $4$, $5$ e $6$,
que estendem o Teorema~1.7 de \cite{ALA20} e explicam o exemplo de Giné et
al.\ citado lá;
(b) um lema elementar de mudanças de sinal, que substitui o uso do teorema de
Briskin--Françoise--Yomdin, unifica todos os casos e mostra que a
conjectura equivale a ``centro $\Rightarrow$ todos os momentos nulos'';
(c) uma fórmula fechada (tipo Cauchy) para o determinante dos momentos, para
qualquer grau;
(d) uma versão ``de hiperplano'' do caso de grau $6$, mais forte que a hipótese
de ortogonalidade, e uma verificação exata assistida por computador da
conjectura de Lijun--Yun para $\deg f\le 7$ e $n\in\{2,4,6\}$
(e, para $n=2$, também com $\deg\fe\le6$ e $\deg\fo\le9$).
Por fim, relacionamos tudo à literatura sobre o problema do centro, momentos e
a conjectura da composição, e listamos problemas em aberto.
\end{abstract}

\tableofcontents

%==================================================================
\section{Contexto e notação}
%==================================================================

Consideramos a equação de Abel
\begin{equation}\label{eq:abel}
 x'(t)=f(t)\,x(t)^3+g(t)\,x(t)^2,\qquad t\in[-1,1],
\end{equation}
com $f,g$ analíticas reais (em geral polinômios). Dizemos que $x=0$ é um
\emph{centro} de \eqref{eq:abel} em $[-1,1]$ se, para todo $\rho$ suficientemente
pequeno, a solução com $x(-1)=\rho$ está definida em $[-1,1]$ e satisfaz
$x(1)=\rho$. Escrevemos
\[
 F(t)=\int_{-1}^t f,\qquad G(t)=\int_{-1}^t g,\qquad
 m_k=\int_{-1}^1 f(t)G(t)^k\,dt\quad(k\ge 0),
\]
e, para $f$ polinomial,
\[
 f=\fe+\fo,\qquad \fe(t)=p(t^2),\qquad \fo(t)=t\,q(t^2).
\]

\paragraph{O problema e a conjectura.}
O problema do centro para \eqref{eq:abel} vem do problema centro-foco de
Poincaré para campos planares: em coordenadas polares e após a transformação
de Cherkas \cite{Che76}, um sistema $\dot x=-y+P_n$, $\dot y=x+Q_n$ com $P_n,Q_n$
homogêneos se torna uma equação de Abel trigonométrica
(ver \cite[\S1.1]{ALA20}). A versão polinomial em um intervalo foi introduzida
por Briskin, Françoise e Yomdin \cite{BFY99,BFY00a,BFY00b} como modelo do problema
clássico. A condição suficiente central é a \emph{condição de composição}
\cite{AL87}: se $F=\tilde F\circ W$, $G=\tilde G\circ W$ com $W(-1)=W(1)$, então
há centro.

\begin{conjLY}[{\cite[Obs.~5.6]{LY01}}]
Seja $f$ um polinômio. A equação $x'=2t\,x^2+f(t)\,x^3$, $t\in[-1,1]$, tem centro
em $x=0$ se, e somente se, $f$ é ímpar.
\end{conjLY}

A versão ``infinitesimal'' (com $\varepsilon f$ no lugar de $f$, para
$|\varepsilon|$ pequeno) foi provada por Lijun--Yun \cite[Teo.~5.5]{LY01} e é caso
particular do teorema de Pakovich sobre o problema do centro paramétrico
\cite{Pak17}. A versão com $\varepsilon=1$ fixo segue em aberto em geral.
Por mudança de escala $x\mapsto\lambda x$, o coeficiente $2t$ pode ser trocado
por $ct$ ($c\ne0$) sem alterar a paridade de $f$; por isso consideramos a
família $g(t)=t^{n-1}$, $n$ par, que contém o caso $n=2$.

%==================================================================
\section{O artigo de Alves--Lemos--Araujo: resumo e observações}
%==================================================================

Os resultados principais de \cite{ALA20} são:

\begin{itemize}[leftmargin=2em]
\item \textbf{Teorema 1.7.} Se \eqref{eq:abel} tem centro, então
 $m_0=m_1=m_2=0$.
\item \textbf{Condição I.} $g$ é ímpar e se anula somente em $t=0$.
 Então $G$ é par, $G(\pm1)=0$, e $G$ é estritamente monótona em $[-1,0]$ e
 em $[0,1]$.
\item \textbf{Teorema 1.10.} Se $g$ satisfaz a Condição I e $p(t^2)$ muda de sinal no
 máximo duas vezes em $[0,1]$, então há centro se, e somente se, $f$ é ímpar.
 A prova usa o Teorema~4.1 de Briskin--Françoise--Yomdin \cite{BFY00a}.
\item \textbf{Teorema 1.12.} Se $g=t^{n-1}$ ($n$ par) e $\deg f\le 5$, então há
 centro se, e somente se, $f$ é ímpar. A prova resolve um sistema $3\times3$
 com determinante explícito
 \[
 \det M=-\frac{16}{15(n+1)(n+3)(n+5)(2n+1)(2n+3)(2n+5)}.
 \]
\item \textbf{Observação 1.8.} Para o exemplo trigonométrico de Giné, Grau e
 Santallusia \cite{GGS14},
 $x'=(\sin\theta-\sin2\theta+\sin3\theta)x^3+(\cos\theta+2\cos2\theta)x^2$,
 há centro, $m_0=m_1=m_2=0$ e $m_3=\pi/2\neq 0$. Logo o Teorema~1.7 não pode
 ser estendido a ``$m_3=0$''.
\end{itemize}

\begin{remark}[Teorema 1.12 é caso do Teorema 1.10]\label{rem:112}
Se $\deg f\le5$, então $p(y)=a_0+a_2y+a_4y^2$ tem no máximo duas raízes, logo
$p(t^2)$ muda de sinal no máximo duas vezes em $[0,1]$. Assim o Teorema~1.12
segue do Teorema~1.10 com $g=t^{n-1}$ (que satisfaz a Condição I). O valor do
Teorema~1.12 está em ter uma prova direta e um determinante explícito; a
Proposição~\ref{prop:cauchy} abaixo dá a fórmula para qualquer grau.
\end{remark}

\begin{remark}[O que de fato limita o grau]
Em todos esses resultados, o argumento tem duas peças:
\begin{enumerate}[label=(\roman*)]
\item \emph{dinâmica}: o centro implica $N+1$ condições sobre os momentos;
\item \emph{álgebra linear/análise}: $N+1$ momentos nulos forçam $p\equiv0$
 sempre que $p$ tem ``complexidade'' $\le N$ (grau ou mudanças de sinal).
\end{enumerate}
A peça (ii) funciona para todo $N$ (Seção~\ref{sec:momentos}). O gargalo é (i):
\cite{ALA20} só obtém $N=2$, e a Observação~1.8 mostra que $m_3=0$ não decorre
do centro em geral. Toda generalização abaixo consiste em extrair mais
informação de (i).
\end{remark}

%==================================================================
\section{\texorpdfstring{Primeira generalização: parte par de grau $\le 4$}{Primeira generalização: parte par de grau <= 4}}
%==================================================================

A nota \emph{A Partial Extension of Theorem~1.12} observa que a prova do Teorema~1.12
usa apenas a parte par de $f$. Reescrevemos o resultado de forma um pouco mais
geral.

\begin{corollary}\label{cor:grau4}
Seja $g$ analítica satisfazendo a Condição~I e seja $f$ um polinômio com
$\deg\fe\le4$ e parte ímpar $\fo$ de grau arbitrário. Então \eqref{eq:abel} tem
centro em $x=0$ se, e somente se, $f$ é ímpar.
\end{corollary}

\begin{proof}
Necessidade: pela Observação~\ref{rem:112}, $p$ é de grau $\le2$, logo muda de
sinal no máximo duas vezes em $(0,1)$, e o Teorema~1.10 (ou o
Teorema~\ref{thm:A}(a) abaixo, que tem prova elementar) dá $p\equiv0$. Para
$g=t^{n-1}$ também vale a prova direta da nota: como $G$ é par, a parte ímpar não
contribui para $m_0,m_1,m_2$, e o sistema $3\times3$ em $(a_0,a_2,a_4)$ tem o
determinante não nulo acima.
Suficiência: se $f$ e $g$ são ímpares, então $f(t)=t\hat f(t^2)$,
$g(t)=t\hat g(t^2)$ e a condição de composição vale com $W(t)=t^2$,
$W(-1)=W(1)$ \cite{AL87}.
\end{proof}

\begin{remark}
Na nota original a hipótese ``$r\ge2$'' é desnecessária: o argumento vale para
qualquer parte ímpar. O ganho real em relação ao Teorema~1.12 é que $\deg f$ pode
ser arbitrário, desde que $\deg\fe\le4$. Como observado, isso já está contido no
Teorema~1.10.
\end{remark}

%==================================================================
\section{Expansão do centro e identidades de ordem superior}\label{sec:identidades}
%==================================================================

Seguimos a mudança de variável usada na nota de grau 6:
\[
 x(t,\rho)=\frac{\rho}{1-\rho\,u(t,\rho)},\qquad u(-1,\rho)=0,
 \qquad u'=g+\frac{\rho f}{1-\rho u}.
\]
Escrevendo $u=\sum_{k\ge0}u_k\rho^k$ e $S:=(1-\rho u)^{-1}=\sum_{k\ge0}S_k\rho^k$,
obtemos a recorrência
\begin{equation}\label{eq:rec}
 u_0=G,\qquad u_k(t)=\int_{-1}^t f\,S_{k-1},\qquad
 S_0=1,\qquad S_k=\sum_{i=0}^{k-1}u_iS_{k-1-i},
\end{equation}
e o centro equivale a $u_k(1)=0$ para todo $k\ge0$. Os primeiros termos são
\[
 S_1=u_0,\quad S_2=u_1+u_0^2,\quad S_3=u_2+2u_0u_1+u_0^3,\quad
 S_4=u_3+2u_0u_2+u_1^2+3u_0^2u_1+u_0^4,
\]
\[
 S_5=u_4+2u_0u_3+2u_1u_2+3u_0^2u_2+3u_0u_1^2+4u_0^3u_1+u_0^5.
\]
Denotemos $H(t)=\int_{-1}^t fG$ e $K_2(t)=\int_{-1}^t fG^2$.

\begin{proposition}[Coeficientes de centro até a ordem 6]\label{prop:ident}
Suponha $G(1)=0$. Então
\begin{align*}
 u_1(1)&=m_0,\qquad u_2(1)=m_1,\qquad u_3(1)=\tfrac12 m_0^2+m_2.
\end{align*}
Se, além disso, $m_0=m_1=m_2=0$, então
\begin{align}
 u_4(1)&=m_3-\frac12\int_{-1}^1 g\,F^2\,dt, \label{eq:u4}\\
 u_5(1)&=m_4-2\int_{-1}^1 g\,G\,F^2\,dt, \label{eq:u5}\\
 u_6(1)&=m_5-\frac12\int_{-1}^1 g\,F^3\,dt-\frac92\int_{-1}^1 g\,G^2F^2\,dt
 +\int_{-1}^1 f\,G^2H\,dt. \label{eq:u6}
\end{align}
\end{proposition}

\begin{proof}
As três primeiras são imediatas: $u_1=F$, $u_2=H$, $u_3=\tfrac12F^2+K_2$. Nas
demais usamos $F(\pm1)=H(\pm1)=K_2(\pm1)=G(\pm1)=0$ (em $-1$ por definição; em
$+1$ por hipótese).

\emph{Ordem 4} (já está na nota de grau 6, com $g=t^{n-1}$):
$u_4(1)=\int f(H+2GF+G^3)$, e $\int fH=-\int fFG$, $\int fFG=-\frac12\int gF^2$.

\emph{Ordem 5}: $u_5(1)=\int fS_4$ com $u_3=\frac12F^2+K_2$. Temos
$\int fF^2=\frac13F^3\big|_{-1}^1=0$, $\int fK_2=-\int fFG^2$,
$2\int fGH=H^2\big|_{-1}^1=0$ e $\int fG^4=m_4$. Logo
$u_5(1)=m_4+2\int fFG^2=m_4+\int G^2(F^2)'=m_4-2\int gGF^2$.

\emph{Ordem 6}: seja $U_4(t)=u_4(t)$, com $U_4'=f(H+2GF+G^3)$. Integrando por
partes, $\int fU_4=-\int fFH-2\int fGF^2-\int fG^3F$. Usando $\int fGK_2=-\int fG^2H$,
$\int fFH=-\frac12\int fGF^2$, $\int fGF^2=-\frac13\int gF^3$ e
$\int fG^3F=-\frac32\int gG^2F^2$, obtemos \eqref{eq:u6}.
\end{proof}

\begin{remark}[Verificação]
As identidades \eqref{eq:u4}--\eqref{eq:u6} foram conferidas simbolicamente
(polinômios genéricos de grau $7$, $g=t^{n-1}$, $n=2,4,6$) e a recorrência
\eqref{eq:rec} foi conferida contra a integração numérica de \eqref{eq:abel}
(script \texttt{codigo/centro\_abel.py}).
\end{remark}

\begin{corollary}[Extensão do Teorema 1.7]\label{cor:17}
Se \eqref{eq:abel} tem centro em $[-1,1]$, então
\[
 m_0=m_1=m_2=0,\qquad m_3=\frac12\int_{-1}^1gF^2,\qquad m_4=2\int_{-1}^1gGF^2,
\]
\[
 m_5=\frac12\int_{-1}^1 gF^3+\frac92\int_{-1}^1 gG^2F^2-\int_{-1}^1 fG^2H .
\]
\end{corollary}

\begin{example}[O exemplo de Giné--Grau--Santallusia]
O argumento da Proposição~\ref{prop:ident} vale em qualquer intervalo $[a,b]$
com $G(a)=G(b)=0$, em particular em $[0,2\pi]$ para equações trigonométricas.
Para o exemplo da Observação~1.8 de \cite{ALA20}, o cálculo numérico (30 dígitos)
dá
\[
 \frac12\int_0^{2\pi} g\,F^2\,d\theta=1{,}5707963267948966192\ldots=\frac\pi2=m_3,
\]
exatamente como prevê o Corolário~\ref{cor:17}. Assim, $m_3\neq0$ nesse exemplo
não é uma anomalia: o quarto coeficiente de centro não é $m_3$, e sim
$m_3-\frac12\int gF^2$.
\end{example}

\begin{remark}[Leitura em termos de integrais iteradas]
$\int gF^2$, $\int gGF^2$, $\int gF^3$ e $\int fG^2H$ são integrais iteradas de
$f\,dt$ e $g\,dt$ (por exemplo $\int fG^2H=\int_{-1<s<t<1}f(s)G(s)f(t)G(t)^2$).
A Proposição~\ref{prop:ident} é, portanto, uma instância explícita da expansão
da aplicação de primeiro retorno em integrais iteradas de Brudnyi \cite{Bru06} e
dos ``momentos duplos'' de Cima--Gasull--Mañosas \cite{CGM13a} e dos
coeficientes de Melnikov de segunda ordem de Briskin--Pakovich--Yomdin
\cite{BPY16}. A novidade aqui é a normalização pela escala $\rho$ (e não por um
parâmetro $\varepsilon$), que faz o momento $m_{k-1}$ aparecer linearmente em
$u_k(1)$, com uma correção não linear em $F$.
\end{remark}

%==================================================================
\section{O papel da paridade}\label{sec:paridade}
%==================================================================

Suponha agora que $g$ satisfaz a Condição~I. Trocando $x$ por $-x$ (o que troca
$g$ por $-g$ e mantém $f$), podemos supor $g>0$ em $(0,1]$. Então $G$ é par,
$G<0$ em $(-1,1)$, e $G|_{[0,1]}$ é um difeomorfismo crescente de $[0,1]$ sobre
$[G(0),0]$; seja $\tau$ sua inversa. Defina, como na nota de grau 6,
\[
 A(t)=\int_0^t\fe,\qquad C(t)=\int_t^1\fo,\qquad
 \ip{\varphi}{\psi}_w=\int_0^1 w\,\varphi\,\psi\,dt .
\]

\begin{lemma}[Decomposição de $F$]\label{lem:par}
Se $m_0=0$, então, para $t\in[0,1]$,
\[
 F(t)=A(t)-C(t),\qquad F(-t)=-A(t)-C(t).
\]
Consequentemente, para todo $j\ge0$,
\[
 \int_{-1}^1 gG^jF^2=-4\ip{A}{C}_{gG^j},\qquad
 \int_{-1}^1 gG^jF^3=2\ip{A^3+3AC^2}{1}_{gG^j}.
\]
Além disso, $m_k=2\int_0^1\fe\,G^k\,dt$ depende só de $\fe$.
\end{lemma}

\begin{proof}
$m_0=0$ e a paridade dão $\int_0^1\fe=0$; as fórmulas para $F(\pm t)$ seguem
como na nota de grau 6. Como $gG^j$ é ímpar,
$\int_{-1}^1 gG^j\Phi(F)=\int_0^1 gG^j[\Phi(F(t))-\Phi(F(-t))]$, e
$(A-C)^2-(A+C)^2=-4AC$, $(A-C)^3+(A+C)^3=2A^3+6AC^2$.
\end{proof}

Combinando com o Corolário~\ref{cor:17}:

\begin{corollary}[Condições de ordem 4 e 5 sob paridade]\label{cor:paridade}
Se $g$ satisfaz a Condição~I e \eqref{eq:abel} tem centro, então
$m_0=m_1=m_2=0$ e
\[
 m_3=-2\ip{A}{C}_{g},\qquad m_4=-8\ip{A}{C}_{gG}.
\]
Na variável $s=G(t)\in[G(0),0]$, temos $g\,dt=ds$, e as condições ficam
\[
 m_3=-2\int_{G(0)}^0(AC)(\tau(s))\,ds,\qquad
 m_4=-8\int_{G(0)}^0(AC)(\tau(s))\,s\,ds .
\]
\end{corollary}

\begin{remark}[Interpretação espectral]
Para $g=t^{n-1}$, a nota de grau 6 interpreta $\ip{A}{C}_g$ via os polinômios de
Jacobi $P_k^{(0,n-1)}(2t-1)$, ortogonais em $L^2((0,1),t^{n-1}dt)$. O
Corolário~\ref{cor:paridade} sugere uma leitura mais intrínseca: na variável
$s=G(t)$ (para $g=t^{n-1}$, $s=(t^n-1)/n$) os pesos $g$ e $gG$ viram $ds$ e
$s\,ds$, isto é, os produtos internos são os \emph{momentos de ordem $0$ e $1$
de $AC$ na variável $G$}, exatamente a mesma variável em que $m_k$ são momentos
de $\fe$. Isso sugere procurar, em ordem $k$, uma relação do tipo
``$m_{k-1}=$ momento de ordem $k-4$ de uma expressão em $A,C$'' mais termos de
ordem superior em $F$; a fórmula \eqref{eq:u6} mostra que, a partir da ordem 6,
aparecem termos cúbicos ($A^3$, $AC^2$) e momentos duplos.
\end{remark}

%==================================================================
\section{A parte \emph{linear}: momentos e mudanças de sinal}\label{sec:momentos}
%==================================================================

\begin{lemma}[Mudanças de sinal]\label{lem:sinal}
Seja $\varphi\in L^1(\alpha,\beta)$ com no máximo $N$ mudanças de sinal e
$\int_\alpha^\beta\varphi(s)s^k\,ds=0$ para $k=0,\dots,N$. Então $\varphi=0$ q.t.p.
\end{lemma}

\begin{proof}
Se $\varphi$ muda de sinal nos pontos $s_1<\dots<s_r$, $r\le N$, então
$\varphi(s)\prod_{i}(s-s_i)$ tem sinal constante e integral nula (pois o produto é
um polinômio de grau $\le N$). Logo é nula q.t.p., e portanto $\varphi=0$ q.t.p.
(Este é o argumento clássico de sistemas de Chebyshev; ver \cite{Kar68,PS72}.)
\end{proof}

\begin{proposition}[Redução a momentos de uma função de uma variável]\label{prop:redu}
Se $g$ satisfaz a Condição~I e $g>0$ em $(0,1]$, então
\[
 m_k=2\int_{G(0)}^0\varphi(s)\,s^k\,ds,\qquad
 \varphi(s)=\frac{p(\tau(s)^2)}{g(\tau(s))},
\]
com $\varphi\in L^1$ e com as mesmas mudanças de sinal de $p$ em $(0,1)$.
Em particular, se $p$ muda de sinal no máximo $N$ vezes em $(0,1)$ e
$m_0=\dots=m_N=0$, então $p\equiv0$.
\end{proposition}

\begin{proof}
Mude a variável $s=G(t)$ em $m_k=2\int_0^1p(t^2)G(t)^k\,dt$; note que
$\int|\varphi|\,ds=\int_0^1|p(t^2)|\,dt<\infty$ (a singularidade de $1/g$ em
$t=0$ é compensada por $ds=g\,dt$). Aplique o Lema~\ref{lem:sinal}.
\end{proof}

\begin{remark}
A Proposição~\ref{prop:redu} é exatamente o que \cite{ALA20} obtém via
\cite[Teo.~4.1]{BFY00a} (a função $\psi$ de \cite[\S3]{ALA20} é $\varphi$ a menos de
um fator), mas com prova de três linhas. Também substitui o argumento com a
identidade de Andréief da nota de grau 6: se $Ma=0$ para a matriz de momentos,
o polinômio $p$ correspondente tem grau $\le N$ em $y=t^2$, logo no máximo $N$
mudanças de sinal, logo $p=0$.
\end{remark}

\begin{corollary}[Reformulação da conjectura]\label{cor:reform}
Se $g$ satisfaz a Condição~I e $f$ é contínua, então $m_k=0$ para todo $k\ge0$
se, e somente se, $f$ é ímpar. Consequentemente, a conjectura de Lijun--Yun (para
$g=t^{n-1}$ ou, mais geralmente, para $g$ com a Condição~I) é equivalente a:
\emph{se \eqref{eq:abel} tem centro, então todos os momentos $m_k$ se anulam}
(implicação (5)/(6) do diagrama de \cite[\S1.3]{ALA20}).
\end{corollary}

\begin{proof}
Se $f$ é ímpar, $fG^k$ é ímpar e $m_k=0$. Reciprocamente, pela
Proposição~\ref{prop:redu}, todos os momentos de $\varphi\in L^1$ no intervalo
compacto $[G(0),0]$ se anulam, logo $\varphi=0$ q.t.p.\ (Weierstrass), isto é,
$\fe\equiv0$. A outra direção da equivalência usa que $f$ ímpar implica centro
(condição de composição com $W=t^2$).
\end{proof}

\begin{remark}
Para $g=t^{n-1}$ isso também segue de Pakovich--Muzychuk \cite{PM09}: os fatores
de composição à direita de $G=(t^n-1)/n$ são, a menos de transformações afins,
$t^d$ com $d\mid n$, e a condição $W(-1)=W(1)$ exige $d$ par, o que força $F$
par. A prova acima é elementar e vale para $f$ apenas contínua (compare com o
teorema de Müntz: $\{t^{nk}\}_{k\ge0}$ é completo porque $\sum 1/(nk)=\infty$
\cite{BE95}).
\end{remark}

Para $g=t^{n-1}$ obtemos uma fórmula fechada para o determinante, generalizando
o do Teorema~1.12.

\begin{proposition}[Determinante tipo Cauchy]\label{prop:cauchy}
Seja $n\ge2$ par, $G(t)=(t^n-1)/n$ e
$M^{(N)}=\bigl(\int_{-1}^1t^{2j}G(t)^k\,dt\bigr)_{k,j=0}^N$. Então
\[
 \det M^{(N)}=\frac{2^{N+1}}{n^{N(N+1)/2}}\;
 \frac{\displaystyle\prod_{0\le i<k\le N}2n\,(k-i)^2}
 {\displaystyle\prod_{i,j=0}^{N}(ni+2j+1)}>0 .
\]
\end{proposition}

\begin{proof}
Expandindo $(t^n-1)^k$, a linha $k$ de $M^{(N)}$ é $n^{-k}$ vezes a linha $k$ de
$C=(2/(ni+2j+1))_{i,j}$ mais uma combinação das linhas anteriores (matriz
triangular unipotente). Logo $\det M^{(N)}=n^{-N(N+1)/2}\det C$, e $C$ é $2$ vezes
uma matriz de Cauchy $\bigl(1/(x_i+y_j)\bigr)$ com $x_i=ni+1$, $y_j=2j$, cujo
determinante é $\prod_{i<k}(x_k-x_i)(y_k-y_i)/\prod_{i,j}(x_i+y_j)$.
\end{proof}

\begin{remark}
A fórmula foi conferida simbolicamente para $n\in\{2,4,6,8\}$ e $N\le4$. Para
$N=2$ ela coincide, a menos das mudanças de escala das linhas feitas em
\cite[eq.~(21)]{ALA20}, com o determinante do Teorema~1.12. A prova também mostra
que $m_0=\dots=m_N=0$ equivale a $\fe\perp t^{nk}$, $k=0,\dots,N$, em
$L^2(-1,1)$: um problema de ortogonalidade de \emph{Müntz} \cite{BE95}. Para
$n=2$ isso significa ortogonalidade a todos os polinômios pares de grau $\le2N$;
por exemplo, para $\deg\fe\le6$ e $N=2$, $\fe$ é múltiplo do polinômio de
Legendre $P_6$.
\end{remark}

%==================================================================
\section{Teoremas unificados}\label{sec:unificados}
%==================================================================

\begin{theorem}\label{thm:A}
Seja $g$ analítica satisfazendo a Condição~I e $f(t)=p(t^2)+tq(t^2)$
polinomial. Em cada um dos casos abaixo, \eqref{eq:abel} tem centro em $x=0$ se,
e somente se, $f$ é ímpar:
\begin{enumerate}[label=(\alph*)]
\item $p$ muda de sinal no máximo $2$ vezes em $(0,1)$
 \emph{(Teorema~1.10 de \cite{ALA20}; contém o Teorema~1.12 e o
 Corolário~\ref{cor:grau4})};
\item $p$ muda de sinal no máximo $3$ vezes em $(0,1)$ e $\ip{A}{C}_g=0$
 \emph{(contém o teorema da nota de grau 6, onde $\deg\fe\le6$ e $g=t^{n-1}$)};
\item $p$ muda de sinal no máximo $4$ vezes em $(0,1)$ e
 $\ip{A}{C}_g=\ip{A}{C}_{gG}=0$ \emph{(em particular, $\deg\fe\le8$)}.
\end{enumerate}
\end{theorem}

\begin{proof}
Necessidade: pelo Corolário~\ref{cor:paridade}, o centro e as hipóteses de
ortogonalidade dão $m_0=\dots=m_N=0$ com $N=2,3,4$ respectivamente; aplique a
Proposição~\ref{prop:redu}. Suficiência: condição de composição com $W=t^2$
\cite{AL87}, como no Corolário~\ref{cor:grau4}.
\end{proof}

A hipótese de ortogonalidade em (b) é implícita, pois $A$ depende de $\fe$, que é
justamente o que se quer provar nulo. Ela pode ser trocada por uma condição
explícita sobre a parte ímpar apenas.

\begin{theorem}[Versão de hiperplano, $\deg\fe\le6$]\label{thm:B}
Seja $g$ como no Teorema~\ref{thm:A}, com $g>0$ em $(0,1]$. Existe um único (a
menos de escala) polinômio par $e\neq0$ de grau $\le6$ com
$\int_{-1}^1eG^k=0$, $k=0,1,2$; além disso,
$\mu:=\int_{-1}^1eG^3\neq0$. Seja $\nu:=\int_{-1}^1eG^4$ e $A_e(t)=\int_0^te$.
Se $f$ é polinomial com $\deg\fe\le6$ e \eqref{eq:abel} tem centro, então $f$ é
ímpar ou $\fe=a\,e$ com $a\ne0$ e
\begin{equation}\label{eq:hiper}
 \ip{A_e}{C}_g=-\frac{\mu}{2}
 \qquad\text{e}\qquad
 \ip{A_e}{C}_{gG}=-\frac{\nu}{8}.
\end{equation}
Em particular, se a parte ímpar de $f$ viola alguma das igualdades em
\eqref{eq:hiper}, então há centro se, e somente se, $f$ é ímpar.
\end{theorem}

\begin{proof}
Pela Proposição~\ref{prop:redu} com $N=3$, a aplicação
$p\mapsto(m_0,\dots,m_3)$ é injetiva no espaço (de dimensão 4) dos polinômios
pares de grau $\le6$. Isso dá a existência e unicidade de $e$ e $\mu\ne0$. Se há
centro, $m_0=m_1=m_2=0$, logo $\fe=ae$ e $A=aA_e$. O Corolário~\ref{cor:paridade}
dá $a\mu=-2a\ip{A_e}{C}_g$ e $a\nu=-8a\ip{A_e}{C}_{gG}$.
\end{proof}

\begin{remark}
Como $C$ depende linearmente de $\fo$, o conjunto excepcional \eqref{eq:hiper} é
a interseção de dois hiperplanos afins no espaço das partes ímpares. A hipótese
de ortogonalidade da nota de grau 6 (que equivale a $\ip{A_e}{C}_g=0$ quando
$a\ne0$) evita o primeiro hiperplano, pois $\mu\ne0$; o Teorema~\ref{thm:B}
mostra que basta evitar \emph{um} dos dois hiperplanos.
\end{remark}

\begin{example}[$n=2$ e $n=4$, $\deg f\le7$]\label{ex:hiper}
Escreva $\fo=b_0t+b_1t^3+b_2t^5+b_3t^7$. Para $g=t$ ($n=2$),
\[
 e(t)=231t^6-315t^4+105t^2-5\;(=16P_6(t)),
\]
e, com $\fe=a\,e/231$, os coeficientes de centro de ordem 4 e 5 são
\begin{align*}
 u_4(1)&=\tfrac{64a}{530675145}\,(255b_1+238b_2+189b_3+765),\\
 u_5(1)&=\tfrac{256a}{10082827755}\,(4845b_0+969b_1+133b_2-63b_3-3876).
\end{align*}
Para $g=t^3$ ($n=4$), $e(t)=385t^6-585t^4+231t^2-15$ e
\begin{align*}
 u_4(1)&=-\tfrac{64a}{5601570975}\,(1615b_0-399b_2-462b_3-693),\\
 u_5(1)&=\tfrac{256a}{128836132425}\,(16169b_0+5313b_1+1817b_2+561b_3-2376).
\end{align*}
Em ambos os casos os dois hiperplanos são independentes, e o conjunto
excepcional é um plano afim de dimensão 2 em $\R^4$.
\end{example}

%==================================================================
\section{Evidência computacional para a conjectura de Lijun--Yun}\label{sec:comp}
%==================================================================

Para ir além do Teorema~\ref{thm:B} é preciso usar $u_6(1),u_7(1),\dots$, que já
são não lineares (ver \eqref{eq:u6}). Como $f$ e $g$ são polinomiais, a
recorrência \eqref{eq:rec} produz $u_k(1)$ como polinômios exatos (coeficientes
racionais) nos coeficientes de $f$. Calculamos esses polinômios e usamos bases de
Gröbner sobre $\Q$ para decidir se o sistema
\[
 u_1(1)=\dots=u_K(1)=0,\qquad a_jz-1=0
\]
(onde $a_j$ é um coeficiente de $\fe$ e $z$ uma variável auxiliar) tem solução
complexa. A base de Gröbner reduzida igual a $\{1\}$ significa que não há
solução, isto é, $u_1(1)=\dots=u_K(1)=0$ força $a_j=0$.

\begin{proposition}[Assistida por computador]\label{prop:comp}
Seja $g=t^{n-1}$. Para $n\in\{2,4,6\}$ e todo polinômio $f$ de grau $\le7$, as
condições $u_1(1)=\dots=u_9(1)=0$ implicam que $f$ é ímpar. Portanto a
conjectura de Lijun--Yun (e sua versão com $t^{n-1}$) vale para $\deg f\le 7$,
$n\in\{2,4,6\}$, \emph{sem} hipótese de ortogonalidade. Para $n=2$, as
condições $u_1(1)=\dots=u_{10}(1)=0$ forçam $f$ ímpar também quando $\deg\fe\le6$ e
$\deg\fo\le9$; logo a conjectura de Lijun--Yun original vale para $\deg f\le9$
com $\deg\fe\le6$ (em particular, para todo $f$ de grau $\le7$).
\end{proposition}

\begin{table}[h]
\centering
\begin{tabular}{@{}cccccc@{}}
\toprule
$n$ & $\deg\fe\le$ & $\deg\fo\le$ & $K$ & conclusão ($u_1(1)=\dots=u_K(1)=0\Rightarrow\fe=0$) \\
\midrule
2 & 6 & 7 & 8 & não (há soluções complexas com $\fe\ne0$) \\
2 & 6 & 7 & 9 & sim \\
4 & 6 & 7 & 8 & não \\
4 & 6 & 7 & 9 & sim \\
6 & 6 & 7 & 9 & sim \\
2 & 6 & 9 & 10 & sim \\
2 & 4 & 7 & 9 & sim (já garantido pelo Corolário~\ref{cor:grau4}) \\
\bottomrule
\end{tabular}
\caption{Testes exatos (aritmética racional, bases de Gröbner em
\texttt{sympy}). Script: \texttt{codigo/centro\_abel.py groebner n E R K}, com
$\deg\fe\le2E$, $\deg\fo\le2R+1$.}
\label{tab:comp}
\end{table}

\begin{remark}
(1) O cálculo é exato e decide a existência de soluções \emph{complexas}; a
conclusão vale, portanto, a fortiori para coeficientes reais.
(2) O fato de $K=8$ não bastar mostra que as condições de ordem 6--8 sozinhas
não excluem o plano excepcional do Exemplo~\ref{ex:hiper}; é a ordem 9 que fecha o
argumento.
(3) O método não dá prova para grau arbitrário: aumentar $\deg\fo$ aumenta o
número de parâmetros livres e, possivelmente, o $K$ necessário.
\end{remark}

%==================================================================
\section{Conexões com a literatura atual}\label{sec:lit}
%==================================================================

\paragraph{Composição e momentos.}
A condição de composição é suficiente para centro \cite{AL87} e foi conjecturada
necessária \cite{BFY99,BFY00a,BFY00b,Chr00,Alw09}. Para o problema
\emph{infinitesimal} (primeira ordem em $\varepsilon$), as condições de centro
são os momentos $\int fG^k=0$ \cite{BFY00a}; Pakovich mostrou que eles não
implicam composição em geral \cite{Pak02}, e Pakovich--Muzychuk resolveram
completamente o problema dos momentos polinomiais: os momentos se anulam se, e
somente se, $F$ é soma de primitivas que satisfazem composições com $G$
\cite{PM09}. Para $g$ com a Condição~I, porém, o anulamento de todos os
momentos já equivale a $f$ ímpar (Corolário~\ref{cor:reform}), o que está por
trás da ``rigidez'' do caso $g=t^{n-1}$: a conjectura de Lijun--Yun é exatamente a
afirmação de que o centro força o anulamento de todos os momentos.

\paragraph{Centros paramétricos e no infinito.}
Briskin--Roytvarf--Yomdin mostraram que as condições de centro ``no infinito''
se reduzem aos momentos \cite{BRY10}. Pakovich resolveu o problema do centro
paramétrico: $x'=gx^2+\varepsilon fx^3$ tem centro para todo $\varepsilon$ se, e
somente se, vale a condição de composição \cite{Pak17}; daí segue a versão
infinitesimal da conjectura de Lijun--Yun. Centros persistentes/universais são
tratados em \cite{CGM12,GGL13}. A conjectura com $\varepsilon=1$ fixo é de
natureza diferente: o Teorema~1.7 e a Proposição~\ref{prop:ident} mostram que,
nessa normalização, o centro só fornece os momentos $m_0,m_1,m_2$ ``puros''; a
partir de $m_3$ eles vêm acoplados a termos não lineares.

\paragraph{Contraexemplos.}
A conjectura da composição é falsa para equações trigonométricas
\cite{Alw09,GGS14} e, como mostraram Giné--Grau--Santallusia, também para
equações polinomiais \cite{GGS19}. Isso torna a conjectura de Lijun--Yun um caso
teste interessante: ela afirma que, para $g=t^{n-1}$, a composição \emph{é}
necessária. Resultados como o Teorema~\ref{thm:A} e a
Proposição~\ref{prop:comp} identificam classes nas quais isso pode ser provado.

\paragraph{Segunda ordem, momentos duplos e geometria do conjunto de centros.}
As integrais da Proposição~\ref{prop:ident} são integrais iteradas no sentido de
Brudnyi \cite{Bru06}. Cima--Gasull--Mañosas deram uma cota explícita para o
número de \emph{momentos duplos} cujo anulamento força composição
\cite{CGM13a,CGM13b}; Briskin--Pakovich--Yomdin estudaram os coeficientes de
Melnikov de segunda ordem e a geometria algébrica do conjunto de centros
\cite{BPY16}; Gavrilov descreveu componentes irredutíveis do conjunto de centros
via teoria de Picard--Lefschetz das funções de Melnikov \cite{Gav20}; e
Álvarez--Bravo--Christopher--Mardešić provaram a conjectura da composição para o
problema infinitesimal em ciclos zero \cite{ABCM21}. O cálculo de
$u_6(1),\dots,u_9(1)$ da Seção~\ref{sec:comp} é uma versão concreta, num caso
específico, dessa abordagem ``de ordem superior''.

\paragraph{Estado da conjectura.}
Até onde pudemos verificar, além de \cite{ALA20} (e de sua versão em arXiv,
1708.01512), não encontramos na literatura avanços posteriores específicos sobre
a conjectura de Lijun--Yun com $\varepsilon=1$. Esta afirmação deve ser
reconferida (MathSciNet/zbMATH, citações de \cite{ALA20}) antes de qualquer
submissão.

%==================================================================
\section{Problemas em aberto e roteiro}\label{sec:abertos}
%==================================================================

\begin{problem}[Grau $6$ completo]
Provar que, para $\deg\fe\le6$ e $\fo$ arbitrário, o plano excepcional
\eqref{eq:hiper} é excluído pelas condições de ordem $\ge6$. A
Proposição~\ref{prop:comp} resolve $\deg\fo\le7$ (e $\deg\fo\le9$ para $n=2$);
nesses testes bastaram $K=9$ ($\deg\fo\le7$) e $K=10$ ($\deg\fo\le9$) condições. Uma prova geral provavelmente exige entender o sinal de
$u_6(1)$ restrito a \eqref{eq:hiper}, usando a decomposição
$\int gF^3=2\ip{A^3+3AC^2}{1}_g$ do Lema~\ref{lem:par}.
\end{problem}

\begin{problem}[Estrutura geral de $u_k(1)$]
Encontrar uma fórmula fechada para $u_{k+1}(1)$, módulo as condições anteriores,
da forma $m_k+\mathcal Q_k(F,G)$, com $\mathcal Q_k$ combinação de integrais
iteradas de grau $\ge2$ em $f$. A Proposição~\ref{prop:ident} dá
$\mathcal Q_3,\mathcal Q_4,\mathcal Q_5$. Uma fórmula geral daria uma
generalização natural do Teorema~1.7 para todos os $k$, e, via
Lema~\ref{lem:par}, condições explícitas em $A$ e $C$.
\end{problem}

\begin{problem}[Ortogonalidades múltiplas]
O Teorema~\ref{thm:A} sugere a família: se $p$ muda de sinal no máximo $N$ vezes
e $\fo$ satisfaz $N-2$ condições de ortogonalidade (nos pesos naturais que
aparecem em $u_4(1),\dots,u_{N+1}(1)$), então centro $\Leftrightarrow$ $f$ ímpar.
Para $N\ge5$ é preciso lidar com os termos cúbicos de \eqref{eq:u6}.
\end{problem}

\begin{problem}[Outros $g$ e o caso trigonométrico]
Os resultados das Seções~\ref{sec:paridade}--\ref{sec:unificados} valem para
qualquer $g$ com a Condição~I. Seria natural (i) repetir a
Proposição~\ref{prop:comp} para $g=t^{n-1}r(t^2)$ com $r>0$; (ii) formular o
análogo trigonométrico (simetria $\theta\mapsto-\theta$ no lugar de $t\mapsto-t$),
onde o exemplo de Giné et al.\ mostra que a correção $-\frac12\int gF^2$ é
essencial.
\end{problem}

\begin{problem}[Análise de $\deg\fe\le 2N$ via Müntz]
Pela Proposição~\ref{prop:cauchy}, $m_0=\dots=m_2=0$ significa
$\fe\perp\{1,t^n,t^{2n}\}$. Descrever explicitamente o núcleo (polinômios de
Müntz--Legendre \cite{BE95}) pode simplificar as contas de $\mu,\nu$ e dos
hiperplanos \eqref{eq:hiper} para todo $n$.
\end{problem}

%==================================================================
\section*{Apêndice: correspondência com os arquivos de partida}
%==================================================================

\begin{itemize}[leftmargin=2em]
\item \texttt{main\_3.tex} (``A Partial Extension of Theorem 1.12''):
 corresponde ao Corolário~\ref{cor:grau4}. O cálculo do sistema linear e do
 determinante está correto (conferido simbolicamente); a Proposição
 \ref{prop:cauchy} dá a versão para qualquer grau.
\item \texttt{abel\_degree6\_orthogonality\_theorem\_1.tex}: o teorema
 corresponde ao Teorema~\ref{thm:A}(b) com $g=t^{n-1}$; a identidade
 \eqref{eq:u4} (Passo 2 da nota) foi conferida simbolicamente. O passo de
 Andréief pode ser trocado pelo Lema~\ref{lem:sinal}; a hipótese de
 ortogonalidade pode ser enfraquecida como no Teorema~\ref{thm:B}; e, para
 $\deg f\le7$, ela é desnecessária (Proposição~\ref{prop:comp}).
\item O Passo 5 da nota de grau 6 (suficiência) pode ser citado diretamente da
 Proposição~1.4 de \cite{ALA20} (Alwash--Lloyd), com $\sigma(t)=t^2$.
\end{itemize}

%==================================================================
\begin{thebibliography}{99}
%==================================================================

\bibitem{ABCM21}
A.~Álvarez, J.~L.~Bravo, C.~Christopher, P.~Mardešić,
Infinitesimal center problem on zero cycles and the composition conjecture,
\emph{Funct. Anal. Appl.} \textbf{55} (2021), no.~4, 257--271.
arXiv:2006.07600.

\bibitem{ALA20}
A.~M.~Alves, A.~Lemos, A.~L.~A.~de Araujo,
A partial positive answer to a Lijun--Yun conjecture,
\emph{Boll. Unione Mat. Ital.} \textbf{13} (2020), 49--59.
doi:10.1007/s40574-019-00203-x. arXiv:1708.01512.

\bibitem{Alw09}
M.~A.~M.~Alwash,
The composition conjecture for Abel equation,
\emph{Expo. Math.} \textbf{27} (2009), 241--250.

\bibitem{AL87}
M.~A.~M.~Alwash, N.~G.~Lloyd,
Non-autonomous equations related to polynomial two-dimensional systems,
\emph{Proc. Roy. Soc. Edinburgh} \textbf{105A} (1987), 129--152.

\bibitem{BE95}
P.~Borwein, T.~Erdélyi,
\emph{Polynomials and Polynomial Inequalities},
Graduate Texts in Mathematics 161, Springer, 1995.

\bibitem{BFY99}
M.~Briskin, J.-P.~Françoise, Y.~Yomdin,
Center conditions, compositions of polynomials and moments on algebraic curves,
\emph{Ergodic Theory Dynam. Systems} \textbf{19} (1999), 1201--1220.

\bibitem{BFY00a}
M.~Briskin, J.-P.~Françoise, Y.~Yomdin,
Center conditions II: parametric and model center problems,
\emph{Israel J. Math.} \textbf{118} (2000), 61--82.

\bibitem{BFY00b}
M.~Briskin, J.-P.~Françoise, Y.~Yomdin,
Center conditions III: parametric and model center problems,
\emph{Israel J. Math.} \textbf{118} (2000), 83--108.

\bibitem{BPY16}
M.~Briskin, F.~Pakovich, Y.~Yomdin,
Algebraic geometry of the center-focus problem for Abel differential equations,
\emph{Ergodic Theory Dynam. Systems} \textbf{36} (2016), 714--744.
arXiv:1211.1296.

\bibitem{BRY10}
M.~Briskin, N.~Roytvarf, Y.~Yomdin,
Center conditions at infinity for Abel differential equations,
\emph{Ann. of Math.} \textbf{172} (2010), 437--483.

\bibitem{Bru06}
A.~Brudnyi,
On the center problem for ordinary differential equations,
\emph{Amer. J. Math.} \textbf{128} (2006), 419--451.

\bibitem{Che76}
L.~A.~Cherkas,
Number of limit cycles of an autonomous second-order system,
\emph{Differ. Equ.} \textbf{5} (1976), 666--668.

\bibitem{Chr00}
C.~Christopher,
Abel equations: composition conjectures and the model problem,
\emph{Bull. London Math. Soc.} \textbf{32} (2000), 332--338.

\bibitem{CGM12}
A.~Cima, A.~Gasull, F.~Mañosas,
Centers for trigonometric Abel equations,
\emph{Qual. Theory Dyn. Syst.} \textbf{11} (2012), 19--37.

\bibitem{CGM13a}
A.~Cima, A.~Gasull, F.~Mañosas,
An explicit bound of the number of vanishing double moments forcing composition,
\emph{J. Differential Equations} \textbf{255} (2013), 339--350.

\bibitem{CGM13b}
A.~Cima, A.~Gasull, F.~Mañosas,
A simple solution of some composition conjectures for Abel equations,
\emph{J. Math. Anal. Appl.} \textbf{398} (2013), 477--486.

\bibitem{Gav20}
L.~Gavrilov,
On the center-focus problem for the equation
$\frac{dy}{dx}+\sum_{i=1}^n a_i(x)y^i=0$, $0\le x\le1$, where $a_i$ are polynomials,
\emph{Ann. Henri Lebesgue} \textbf{3} (2020), 615--648.

\bibitem{GGL13}
J.~Giné, M.~Grau, J.~Llibre,
Universal centres and composition conditions,
\emph{Proc. London Math. Soc.} \textbf{106} (2013), 481--507.

\bibitem{GGS14}
J.~Giné, M.~Grau, X.~Santallusia,
Universal centers in the cubic trigonometric Abel equation,
\emph{Electron. J. Qual. Theory Differ. Equ.} (2014), no.~1, 1--7.

\bibitem{GGS19}
J.~Giné, M.~Grau, X.~Santallusia,
A counterexample to the composition condition conjecture for polynomial Abel
differential equations,
\emph{Ergodic Theory Dynam. Systems} \textbf{39} (2019).
doi:10.1017/etds.2018.16.

\bibitem{Kar68}
S.~Karlin,
\emph{Total Positivity}, Vol.~I, Stanford University Press, 1968.

\bibitem{LY01}
Y.~Lijun, T.~Yun,
Some new results on Abel equations,
\emph{J. Math. Anal. Appl.} \textbf{261} (2001), 100--112.

\bibitem{Pak02}
F.~Pakovich,
A counterexample to the ``composition conjecture'',
\emph{Proc. Amer. Math. Soc.} \textbf{130} (2002), 3747--3749.

\bibitem{PM09}
F.~Pakovich, M.~Muzychuk,
Solution of the polynomial moment problem,
\emph{Proc. London Math. Soc.} \textbf{99} (2009), 633--657.

\bibitem{Pak17}
F.~Pakovich,
Solution of the parametric center problem for the Abel differential equation,
\emph{J. Eur. Math. Soc.} \textbf{19} (2017), 2343--2369.

\bibitem{PS72}
G.~Pólya, G.~Szegő,
\emph{Problems and Theorems in Analysis}, Vol.~II, Springer, 1976.

\end{thebibliography}

\end{document}
